\documentclass[quick,con]{debate}
\motion{我厌蠢，我错了 / 我厌蠢，我没错}
\stance{反方}{我没错}
\begin{document}
\debatetitle
\begin{thesisbox}常人的厌蠢，平常冲着教过了还犯，看得有据，怪得公道；恼里那句「你本可以」，是对能负责的人提要求，这回放下，下回照样说、照样听。问完三句，这份厌平常的样子站得住。\end{thesisbox}

\section{定义与判准（标准表述）}
\begin{fields}
\field{厌蠢}{厌蠢，是觉得别人犯蠢时生出的嫌恶、反感，常带着想躲开的心；它可以只冲着那件事，也可以冲着这个人、看低他；它是一种态度，不等于骂人或欺负人。}
\field{蠢}{蠢，是在一般人看来本该想得到、做得对的事上，明显想错、做错；可能出于不会，也可能出于不肯，原因不在定义里预设。\sub{被问才答}{一般人，是了解他处境、讲道理的一般人；我，是一个普通的、会说出这句话的我。}}
\field{判准}{判断这份厌对不对，就看一个清醒的我问完三句：看准了吗？对他公道吗？划得来吗？问完觉得这份厌平常的样子就该改，就是错了；站得住，只在个别场合调调分寸，就是没错。\sub{门槛}{要改的如果是这份厌平常的样子，就是这份厌本身该改，哪怕改完还剩一点反感；要改的如果只在个别场合冒出来，那叫调分寸。平常就这样、只是改得掉，照样算错了。}\sub{义务}{我方：问完三句，这份厌平常的样子站得住；对方：它平常的样子该改。}\sub{越线}{这回先不劝、一时走开、不再主动提醒，不越线；他开口也不听、合作只下指令不给理由，越线，我们认要调分寸。}}
\field{30 秒口径}{正方用了中立版：「定义与判准，和对方大体一致；我们争的是门槛。」把看低写死在厌里：「汉典里，看不起是鄙视的释义。」把门槛压成有任何毛病：「照这个标准，世上没有不错的情绪。」点完重述中立版。}
\end{fields}

\section{我方论点}
\begin{tabularx}{\linewidth}{@{}L{4.6em}L{8em}Y L{6.4em}@{}}\toprule
\thead{标签}&\thead{一句话}&\thead{核心数据 / 学理 / 案例}&\thead{出处}\\\midrule
\textbf{教过了还犯}（准·归责）&被教过、提醒过，仍一再出错、不问不改；看得有据，怪得公道&元分析《A meta-analytic review of help giving and aggression from an attributional perspective》：64 项研究、12,000 余名被试，该负责就生气、控制不了就同情（只核摘要，对象是求助者与冒犯者）。学理：罪与病。案例：同事工作超过两年常填错表格，多次劝告，一度降至冰点，现在很微妙&Rudolph 等 2004；Weiner 1993；每日商报 2024\\
\textbf{恼他不弃他}（公·尊重）&收回的是这一次的评价，作为人的尊重不动；「你本可以」是对能负责的人提要求&学理：态度完全客观才不能讲理（恼里带着讲理是推论）。数据：《Beat them or ban them》3 项研究，愤怒长远和解、轻蔑推人出去（须认 often occur together、may develop）。案例：孔子斥子路「野哉」后讲完正名，胡氏注子路终不喻。康德德性论第三十九节整句引（自译）&Strawson 1962；Fischer \& Roseman 2007；《论语·子路》；Kant 1797\\
\bottomrule\end{tabularx}
\textbf{其余材料}：Hutcherson \& Gross 2011《The moral emotions》，与自己相关推高愤怒、轻蔑在评判无能行为时起作用（整句引）；Brady 等 2021《How social learning amplifies moral outrage expression…》，推特上道德愤怒被点赞后更常表达；Graham 1984《Communicating sympathy and anger to Black and White children》，儿童样本；Miceli \& Castelfranchi 2018《Contempt and disgust》，轻蔑不看好改善；Brooks 等 2015，求助被看作更有能力；Edmondson 1999，51 个工作团队；Heslin 等 2006《Keen to help?》，3 项研究，不测厌；Roberts 2013 书评《Hard Feelings》概括 Bell。

\section{对方论点 → 一句话反驳}
\begin{tabularx}{\linewidth}{@{}L{6.8em}Y Y@{}}\toprule
\thead{对方会说}&\thead{我方一句话回}&\thead{对方追问时}\\\midrule
平常凭比较，不凭记录&厌跟着代价走，代价落在一起做事的人身上，记录现成；评论区那句我们认：看错了，个别场合&“新京报写同学同事也是‘我都可以’”→ 全认；只取场景，记录由每日商报那一案补\\
会了的人估不准、看不见&这些实验估的都是一个他一无所知的人；同事学过没有，查得到，用不着估&“方向固定就是平常”→ 有记录时尺子不上场；我们不评偏多少\\
劝了还犯，恼会长成轻蔑&可能不是平常，预测因素不是必要条件；还在劝，就是还觉得话管得到他&“常常一起出现”→ 整句认；往哪走，看回头还讲不讲\\
回头不讲理，越过尊重&三句一起读：多次劝告、一度降至冰点、现在很微妙；之后讲不讲理，报道没写&“怎么还是那么蠢”→ 没标说话人；“怎么还是”还在要求他（推论）\\
厌蠢就是轻蔑、傲慢&看低他这件事上的表现可以，拒绝作为人的尊重不行；康德整句引&“心里也会漏出来”→ 管住表现，守住尊重\\
替代态度：先开口问&没有记录先问，本来就在我们的厌里；存在更好的，不等于这个就错了&“先问要学才有”→ 问是在拿记录（推论）；劝本身就是开口\\
\bottomrule\end{tabularx}

\section{必问与必答}
\begin{points}{必问（不答就点名、重复、不换题）}
\item 您方举的例子里，他有没有被教过？（全场问题）\quad 证平常的材料，除了评论区和一则报道，还有吗？
\item 回头不讲理，有直接材料吗，还是推论？\quad 恼火您方不反对，去掉看低剩下的那份，算不算错？
\end{points}
\begin{points}{必答（15 秒正面回答）}
\item 交涉还是走开？→ 可以一时走开，但不把他划出讲理的范围。\quad 只在心里厌算不算？→ 算，错不错看判准。
\item 看准了就能看低他吗？→ 看低他的能力可以，拒绝给他作为人的尊重不行。\quad 对真正的笨？→ 厌了，那一次是看错了。
\end{points}

\section{自由辩战场概要}
\begin{tabularx}{\linewidth}{@{}L{7.2em}Y Y@{}}\toprule
\thead{战场}&\thead{开场问题}&\thead{收束语}\\\midrule
平常在哪（准）&您方举的例子里，他有没有被教过？&厌跟着代价走，那里记录现成\\
恼了之后（尊重）&回头不讲理，有直接材料吗，还是推论？&这回放下，下回照样说；恼他，不弃他\\
替代（门槛）&对劝过多次的同事，先开口问，还能问出新东西吗？&存在更好的，不等于这个就错了\\
\bottomrule\end{tabularx}
前 90 秒战场一（抛完必问），中 90 秒战场二，后 60 秒战场三并收束；一辩开场，四辩收束；不开“划得来吗”的进攻。

\section{如果……就……}
\begin{contingency}
\ifthen{对方拿出没预料到的数据}{先问机构、年份、报告名，再问测的是不是被教过的人；报厌蠢比例的，请出示来源。}
\ifthen{论点一被打成“熟人也凭比较”}{收缩到论点二：有，不等于平常。被问推论：直接认，给依据。}
\end{contingency}

各环节：正一立论 180 秒 → 反二质询 120 秒双边 → 反一立论 180 秒 → 正二质询 120 秒双边 → 反二驳论 120 秒 → 正二驳论 120 秒 → 四辩对辩各 90 秒 → 正三盘问 90 秒 → 反三盘问 90 秒 → 正三小结 120 秒 → 反三小结 120 秒 → 自由辩各 240 秒 → 反四结辩 210 秒 → 正四结辩 210 秒；奇袭（反三质询 / 反一申论）首选自由辩之后。

\begin{recordcard}
\record{正一立论结束}{正方定义与判准}{\opt{正方用了双方共用的中立定义与判准}\opt{正方把看低写死在厌里}\opt{正方的定义与判准不在以上几种时}}{反一立论·定义与判准}
\record{正一立论结束}{对方立论的主干}{\opt{对方立论建立在「厌蠢平常凭比较，不凭记录」上}\opt{对方立论建立在「会了的人估不准、自己看不见」上}\opt{对方立论建立在「劝了还犯，恼会长成轻蔑」上}\opt{对方立论建立在「回头不讲理，越过了尊重」上}\opt{对方立论建立在「厌蠢就是轻蔑、傲慢」上}\opt{对方立论的主干不在以上几条时}}{反二驳论（选 2）}
\record{反二质询结束}{质询拿到了什么}{\opt{质询拿到了对方承认凭记录的不争}\opt{质询拿到了对方承认对那件事的恼火不反对}\opt{质询里对方回避了「他有没有被教过」}\opt{质询没有拿到明确成果}}{反一立论·开头}
\record{反三盘问结束}{链一的结果}{\opt{对方承认平常的材料只有评论区和报道}\opt{正二、正四答得不一样}\opt{对方回避了链一的问题}\opt{链一的结果不在以上几种时}}{反三小结·链一}
\record{反三盘问结束}{链二的结果}{\opt{对方承认回头不讲理是推论}\opt{对方承认下回照样说不算越线}\opt{链二的结果不在以上几种时}}{反三小结·链二}
\record{正三小结结束}{正三引向哪里}{\opt{正三把战场引向「没有记录时凭比较」}\opt{正三把战场引向「怎么还是那么蠢」}\opt{正三把战场引向「先问是要学才有」}\opt{正三小结的方向不在以上几条时}}{反三小结·回应正三}
\record{发动奇袭时}{在哪个窗口}{\opt{在盘问小结之后发动}\opt{在自由辩之后发动}\opt{发动时机不在以上几个窗口时}}{奇袭申论}
\record{自由辩结束}{全场主战场}{\opt{全场主战场落在「平常凭比较还是凭记录」}\opt{全场主战场落在「看准之后还讲不讲理」}\opt{主战场不清楚时}}{反四结辩·归纳分歧}
\record{自由辩结束}{对方最有力的一点}{\opt{对方全场最有力的是「会了的人估不准、自己看不见」}\opt{对方全场最有力的是「轻蔑可能从一次次恼火里长出来」}\opt{对方全场最有力的是「说了那么多遍，怎么还是那么蠢」}\opt{对方最有力的一点不在以上几条时}}{反四结辩·最强点}
\record{自由辩结束}{正方主打的路线}{\opt{正方全场主打「平常就是凭比较」}\opt{正方全场主打「就算看准了，也会滑到看低人」}\opt{正方全场主打「替代做得到，所以该改」}\opt{正方全场路线不清楚时}}{反四结辩·封路}
\end{recordcard}
\end{document}
